\documentclass[12pt]{article}
\usepackage[utf8]{inputenc}

%%% I added a few packages. This next one allows you to set the margins.
\usepackage[top=.5in,bottom=1in,right=.75in,left=.75in]{geometry}
\usepackage{multicol}
%\usepackage{newspaper}
\setlength{\columnsep}{1cm}
\usepackage{blindtext}
\usepackage{varwidth}

%%% This one does boxes in a really nice way.
\usepackage{tcolorbox}
\usepackage{mdframed}
\usepackage{minibox}
\usepackage{graphicx}
\graphicspath{ {./images/} }
\usepackage{epstopdf}
  \epstopdfDeclareGraphicsRule{.pdf}{png}{.png}{convert #1 \OutputFile}
  \DeclareGraphicsExtensions{.png,.pdf}
\usepackage{hyperref}
\hypersetup{
    colorlinks=true,
    linkcolor=blue,
    filecolor=magenta,      
    urlcolor=cyan,
    pdftitle={MSCSMess},
    pdfpagemode=FullScreen,
    }
\urlstyle{same}
 \usepackage{fix-cm}
 \usepackage{soul}
 \usepackage{graphicx}
 \usepackage{multicol}


\newcommand{\ds}{\displaystyle}
\newcommand{\ZZ}{\mathbb{Z}}
\newcommand{\QQ}{\mathbb{Q}}
\newcommand{\RR}{\mathbb{R}}
\newcommand{\CC}{\mathbb{C}}
\newcommand{\FF}{\mathbb{F}}
\newcommand{\xx}{\mathbf{x}}
\newcommand{\uu}{\mathbf{u}}
\newcommand{\vv}{\mathbf{v}}

\newcommand{\bolder}{\textbf}
 
 
\title{\fontsize{60}{70}\selectfont {\bf {MSCS} 
\includegraphics[scale=.3]{images/OleLion.png} {Mess}}}
\author{Department of Mathematics, Statistics, and Computer Science \\ St.\@ Olaf College, Northfield, MN 55057 \vspace{-2ex}} 
\date{\today\, $|$ Volume 54, No.\@ 10 \line(1,0){500}}


\begin{document}

\maketitle \vspace{-0.5cm}

Hey Oles! Hopefully you got to see the Northern Lights last Tuesday, they were spectacular! As it gets darker and colder outside and you're making the final push towards the end of classes, be sure to take mental breaks, destress, and take care of yourself. In the issue below, check out the colloquium happening next week and information on applying to be a TA for the MSCS department!


 \begin{multicols}{2}

 \begin{tcolorbox}[colback=white, colframe=black]
    \begin{center}
    \title{{\large \bolder{Tea Time Event and Math Colloquium 11/17
}}}
    \end{center}
\end{tcolorbox}

On Monday, the Math Club will be hosting a ``crafternoon" during Tea Time in RMS! Immediately following the event, math faculty candidate Katharine Adamyk will be hosting a colloquium on algebraic topology. See the posters below for details!

\begin{tcolorbox}[colback=white, colframe=black]
    \begin{center}
        \title{{\large \bolder{MSCS TA/Grader Applications}}}
    \end{center}
\end{tcolorbox}

The MSCS department is now accepting Grader/TA applications for J-term/Spring 2026. If you are interested in working with us please apply \href{https://forms.gle/Ejy4MdjvExAKWxiM7}{here}. The deadline for applying is Tuesday, November 18th at 11:59pm. Please note the qualifications below:
\begin{itemize} \itemsep -4pt
    \item Completed the course you are applying for with an A or B grade.
    \item Be on-campus (not studying abroad).
    \item Ability to maintain confidentiality.
    \item Complete work in a timely manner.
\end{itemize}

\begin{tcolorbox}[colback=white, colframe=black]
    \begin{center}
    \title{{\large \bolder{MSCS student profile in the St. Olaf Magazine
}}}
    \end{center}
\end{tcolorbox}

There is an excellent article about Jessica Schmidt--math/bio major, SDS/math bio concentrator, class of 2026--in the current issue of the \href{https://wp.stolaf.edu/magazine/}{St. Olaf Magazine}, titled ``The Cipher Protocol." Jessica is a ``master codebreaker" after having been the only individual to have solved all three codebreaking challenges in the annual Kryptos competition. Read the article (and the follow-up article on the MSCS department and Math program) and congratulate Jessica.

\begin{tcolorbox}[colback=white, colframe=black]
    \begin{center}
    \title{{\large \bolder{CSCI 121 Principles of Computer Science (J-term)
}}}
    \end{center}
\end{tcolorbox}

Dive into computer science with Professor Houssou! This introductory, hands-on course will develop your analytical thinking and problem-solving skills through daily programming exercises. You will apply programming concepts to real-world projects across various fields like science and engineering, as well as music, sociology, finance, art, and more. 

The smaller class size during J-term allows you a more personalized experience into the realm of computer science with one-on-one support from Professor Houssou and close collaboration with your classmates. This is also a great opportunity to build up your technology skills in an ever-increasing technical world.

\begin{tcolorbox}[colback=white, colframe=black]
    \begin{center}
    \title{{\large \bolder{ CS Unplugged   }}}
    \end{center}
\end{tcolorbox}

This year the CS program is hosting CS Unplugged: a four-part series on internships, jobs, and life as a computer scientist. These sessions will be informal, student-centered, and are a part of a larger goal to get CS students building community within the program. Fill out the form \href{https://docs.google.com/forms/d/e/1FAIpQLScfAnexFhdd-ZAHY0RDX8MpB6A7D7cFtjFYITs-VHnXKzg8XQ/viewform}{at this link} to give your input on what you want to hear at these events!

The first two events will be held on Wednesday, Nov 12th from 3:30-4:30pm, and Wednesday, Nov 19th from 3:30-4:30pm, both in RNS 203. At the end of each event, prizes will be given out! See the poster below for more info.

\begin{tcolorbox}[colback=white, colframe=black]
    \begin{center}
    \title{{\large \bolder{Problem Solving Group}}}
    \end{center}
\end{tcolorbox}
The Math Problem Solving Group meets weekly to have fun, solve math problems, eat candy, and prepare for upcoming math contests. Meetings are 11:30am-12:30pm, in RNS 207 (the Center for Interdisciplinary Research room). Everyone is welcome, regardless of prior experience in mathematics or math contests.

The first contest will be the Jerry Heuer Team Contest on the morning of Saturday, November 15. It's a fun, accessible contest for teams of up to three students. For more information, or to be added to the problem solving group email list, contact Prof. Wright at wright5@stolaf.edu.


\begin{tcolorbox}[colback=white, colframe=black]
    \begin{center}
    \title{{\large \bolder{DAAD RISE Internships}}}
    \end{center}
\end{tcolorbox}

Are you interested in a summer STEM internship/fellowship program for undergraduates in Germany, requiring no German-language skills? Check out the DAAD Research Internships in Science and Engineering (RISE) website \href{https://www.daad.de/rise/en/}{at this link}. The application window is October 15th - November 30th.

\end{multicols} \vspace{0.5cm}

\vfill

\begin{center}
\includegraphics[width=0.9\textwidth]{images/Colloquium Ian Whitehead.png}
\end{center}

\begin{center}
\includegraphics[width=0.9\textwidth]{images/Colloquium Katharine Adamyk (1).png}
\end{center}

\begin{center}
\includegraphics[width=0.9\textwidth]{images/Math Club Crafternoon.png}
\end{center}

\begin{center}
\includegraphics[width=0.9\textwidth]{images/CSCI 121 J-term Course.png}
\end{center}

\begin{center}
\includegraphics[width=0.9\textwidth]{images/Tea Time Poster.png}
\end{center}

\noindent{To submit an article, event, or anything else for publication in the Mess, email messeditor@stolaf.edu. Visit \href{https://wp.stolaf.edu/mscs/mscs-mess/}{http://wp.stolaf.edu/mscs/mscs-mess/} for a digital archive of previous MSCS Mess issues.}\\
	
\hfill Zach Martin, Editor
	
\hfill Kyle Salois, Faculty Advisor


 

\end{document}